
\newcommand{\bolarep}[1]{\Tcircle{\mbox{\scriptsize #1}}}
\newcommand{\arbrerep}[2]{\pstree[levelsep=25pt,treesep=12pt]{\bolarep{#1}}{#2}}


\Problem{Arbres --- Mateixes fulles}


\Statement

En un arbre binari, una fulla és un node sense cap fill. 

\medskip

\begin{latexonly}
Per exemple, l'arbre
\[
\arbrerep{3}
{
  \pstree{\bolarep{0}}
  {
    \pstree{\bolarep{7}}
    {
      \Tn
      \bolarep{4}
    }
    \pstree{\bolarep{2}}
    {
      \Tn
      \Tn
    }
  }
  \pstree{\bolarep{5}}
  {
    \pstree{\bolarep{4}}
    {
      \Tn
      \Tn
    }
    \pstree{\bolarep{7}}
    {
      \pstree{\bolarep{6}}
      {
        \Tn
        \pstree{\bolarep{1}}
        {
          \Tn
          \Tn
        }
      }
      \Tn
    }
  }
}
\vspace*{-0.5cm}
\]
té com a fulles (d'esquerra a dreta) 4, 2, 4, 1.
\end{latexonly}

\medskip

Feu una funció que, donats dos arbres binaris de naturals, 
digui si tenen o no les mateixes fulles d'esquerra a dreta. 
Resoleu el problema utilitzant una funció que generi 
les fulles d'esquerra a dreta d'un arbre binari. 
Digueu quin és el cost asimptòtic en el cas millor i en el cas pitjor de les vostres funcions.

\Observations

\begin{itemize}
    
    \item Descarregueu-vos el fitxer \texttt{code.py}. El programa principal,
    les estructures de dades, 
    la lectura de l'arbre i l'esquelet de les funcions ja se us dónen implementats.

    \item
    El programa principal serveix per provar les funcions i llegeix $p$ parells d'arbres en preordre (amb valors $-1$ pels arbres buits)
    i escriu, per a cada parell, les fulles d'esquerra a dreta dels dos arbres i si tenen o no les mateixes fulles d'esquerra a dreta. 
    
    \item Us pot anar bé fer servir la funció @zip_longest@ d'@itertools@ que funciona com @zip@ però afegeix @None@s si una de les dues seqüències acaba abans que l'altra.
    
\end{itemize}


\SampleOneCol
